\chapter{实现效果与实验验证}
\section{验证方法与版本归属}
验证分为硬件实现报告、RTL与主机数字检查、真实板端数据、冻结EXE检查和真人交互五类。仿真检查状态转换和边界条件，实板检查接口与数据完整性，真人操作检查交互效果，打包检查确认交付文件完整。

最终正式软件为1.6.0，沿用1.5.0的\code{ec0351f8}固件。三敲与中止的数字和运行证据属于1.5.0；界面重构实测属于1.3.0，板载面板属于1.4.0；最终集成检查说明这些能力在1.6.0中的兼容关系。前期触发策略实验保留其原构建归属，仅用于分析策略权衡。

\section{最终资源与100 MHz时序}
\begin{table}[H]\centering\caption{最终固件ec0351f8的实现资源}\small
\begin{tabular}{lrrr}\toprule 资源 & 使用量 & 器件总量 & 使用率\\\midrule
Slice LUTs &15118&63400&23.85\%\\
Slice Registers &25096&126800&19.79\%\\
Block RAM Tile &119.5&135&88.52\%\\
DSP48E1 &9&240&3.75\%\\\bottomrule\end{tabular}
\end{table}
\fig{resources}{最终固件主要资源占用率}{根据Vivado布线后资源报告生成；BRAM是当前扩展的主要资源约束。}

\begin{table}[H]\centering\caption{最终固件时序摘要}\small
\begin{tabularx}{\linewidth}{p{54mm}Y}\toprule 指标 & 实现结果\\\midrule
系统目标周期/频率 &10.000 ns / 100 MHz\\
建立时间最差裕量WNS &+0.102 ns\\
保持时间最差裕量WHS &+0.026 ns\\
建立/保持总负裕量 &0.000 ns / 0.000 ns\\
建立/保持违例端点 &0 / 0\\
慢速串行输出组最差裕量 &+11.933 ns（20 ns传播上限）\\\bottomrule
\end{tabularx}\end{table}

以上为已布线实现结果，说明该构建满足设定约束。外部异步输入和人眼显示按XDC中的同步器边界与显示性质设置例外；并不宣称所有外部引脚都有同步外设时序模型。实现DRC仍保留建议性警告，例如部分DSP寄存器流水建议，原始报告随源码保存，便于区分时序通过与进一步优化空间。

\section{多传感器、完整事件与界面效果}
历史多传感器实板验证在约412.45秒内取得8110次监测快照，覆盖全部16路开关高低电平和5个方向键按下/释放。该轮板载温度范围23.44～24.09℃，FPGA温度32.13～33.18℃；VCCINT约1.000 V、VCCAUX约1.792 V、VCCBRAM约1.000 V。该结果用于证明通道读回与动态监测，版本为\code{f4e259e1}，不作为最终版重新标定结果。

\shot{figures/monitor.png}{实板监测页面与多源信息布局}{1}{实板面板验收阶段的实际软件截图（1.4.0），1.6.0沿用同一监测界面；最终新增游戏页另列。}

在最终1.6.0集成回归中，退出识曲后成功导出一条24000个PCM点、281组XYZ的完整事件，整段校验、波形窗口及WAV播放API通过。图\ref{fig:event}使用该原始负载，每80点保留极值绘制包络，完整数据另行保存。

\fig{event}{最终版本完整1.5秒事件的声音包络}{横轴相对记录锚点，范围为触发前0.5秒至后1秒；仅为绘图进行极值压缩，不改原始PCM。该回归事件不用于认定新的敲击识别准确率。}

\section{声振融合的前期对照结果}
前期实验共30条实板窗口，开发集和留出集各15条，包含轻敲、纯声音与安静输入。阈值及配对窗由开发集确定后固定；留出集含5次轻敲和10次非敲击。实际数据回送板端处理，并与独立整数参考逐窗比较，共4500个窗口的特征和三种决策一致。以下保留原构建\code{b247b381}的试次级结果。

\begin{table}[H]\centering\caption{原版触发策略留出集结果}\small
\begin{tabular}{lrrrrr}\toprule 策略&TP&FP&FN&召回率&F1\\\midrule
\HoldoutRows
\bottomrule\end{tabular}\end{table}

纯声音策略较敏感，但非敲击输入也产生触发；振动与联合策略在该小样本中均无误报，联合策略的召回低于振动单路。结果表明，融合条件实现了可解释的取舍，同时提示阈值和配对窗仍需适配实际环境。最终三敲属于连续动作入口，不能把这张单次事件表直接当作三敲唤醒的灵敏度。

\section{在线识曲与三敲、中止验证}
保留的三首板载麦克风真实录音，经冻结版识曲工作进程回放均匹配正确。该轮调用在线服务，使用历史录音输入，不是新一轮手机外放采样；表中耗时为工作进程识别耗时，不含20秒录音和串口导出。

\begin{table}[H]\centering\caption{真实板载录音的Shazam回放结果}\small
\begin{tabularx}{\linewidth}{Ylll}\toprule 输入歌曲&来源&结果&识别/s\\\midrule
\ShazamRows
\bottomrule\end{tabularx}\end{table}

Shazam替换后的回放说明，硬件录音能够进入在线识别链路并得到可核对的歌曲标签。早期AudioTag阶段有未匹配和低置信候选，历史数据保留；第三方服务、录音片段与策略变化后的结果分别统计。离线歌单作为备用保留，开发调参数据不外推为正式准确率。

三敲及中止升级已通过主机行为、RTL状态边界、冻结资源和实际运行检查；1.5.0运行记录覆盖真实录音/传输中止与重新就绪，另有300秒安静运行未出现敲击或新录音。最终1.6.0继续使用该固件。当前材料未提供独立真人三敲灵敏度统计，因此不报告误唤醒率或泛化准确率。

\section{体感游戏与最终集成验收}
\shot{figures/game.png}{最终1.6.0主程序中的八关体感游戏}{1}{实际冻结EXE连接真实开发板的截图；进度由已验收独立版复制，图中17/24为首次通关的累计星级。}

真人实板记录覆盖四向校准、BTNC刹车及八关通关，累计读到9次物理中键按下、109条刹车控制记录。各关首次通关结果如下，保留失误及用时，不把数字控制器的结果替代真人体验。

\begin{table}[H]\centering\caption{真人实板八关首次通关记录}\small
\begin{tabular}{rrrrl}\toprule 关卡&关内用时/s&星级&死亡次数&结果\\\midrule
\LevelRows
\bottomrule\end{tabular}\end{table}

最终合并侧重共享设备与原功能回归，没有要求重新重玩八关。100项自动检查通过；真实集成期间只存在一个板端采集线程，游戏退出后精确恢复触发配置；识曲启用时拒绝进入实板游戏；正常关闭与冻结EXE实际校准入口检查通过。

\begin{table}[H]\centering\caption{最终软件集成与交付检查}\small
\begin{tabularx}{\linewidth}{p{39mm}Y}\toprule 检查项&已保存结果\\\midrule
主机自动检查 &100项行为、会话、配置恢复及原功能保护检查通过。\\
原资源一致性 &最终EXE内158项原资源逐字节核对通过，固件保持ec0351f8。\\
隔离启动 &中文路径、独立用户目录、干净PATH下自检通过。\\
真实完整事件 &24000个PCM点、281组XYZ；完整性、波形窗口及播放API通过。\\
共享服务 &一个UART拥有者，识曲互斥、游戏退出配置恢复通过。\\
界面及生命周期 &懒加载入口、双主题、1060像素窄窗口、离页暂停及手动继续通过。\\\bottomrule
\end{tabularx}\end{table}

带真实串口流的10秒键盘控制绘制检查记录602帧，约60.22 FPS，帧间隔P95为16.91 ms；主机解码消息发布到Canvas更新提交的P95约56.44 ms。该项是键盘控制下的短时软件调度测量，不是完整物理动作到显示器的端到端延迟，也不外推为所有电脑性能。

\section{结果解释与适用边界}
项目已形成可运行、可演示且数据可核对的系统。最强的证据是完整采集与事务式导出、已布线时序、最终EXE资源一致性、模式互斥恢复以及真人八关体验。识曲三首回放结果说明已接通可用的语义识别应用，但小规模固定录音不代表任意声源和环境下的准确率。

温度、电压、声压、加速度和倾角未通过外部仪表绝对标定；实时预览可能出现序号缺口，完整数据应以校验后的事件文件为准。另一台电脑、第二块板和常规FTDI实板路径尚无对应实测，本机结果不自动扩展到这些条件。上述边界不妨碍已完成的功能验收，同时为后续完善提供明确方向。
